\paragraph{一般纤维扫描的循环型卷积公式}
设可逆位点诱导图分解为路径分量直并 \(\Gamma_x=\bigsqcup_{j=1}^r P_{\ell_j}\)，从而
\(
F_x\cong \prod_{j=1}^r \mathsf{Ind}(P_{\ell_j})
\)
且扫描元在分量坐标下按坐标逐因子作用（定义 \ref{def:pom-toggle-coxeter-element}）。
定理 \ref{thm:pom-toggle-scan-lcm-tensorization} 给出点周期的 \(\mathrm{lcm}\) 规律；下面给出循环分解层面的精确计数卷积。

\begin{proposition}[循环条数的 \(\mathrm{lcm}\)-卷积闭式]\label{prop:pom-toggle-scan-cycle-convolution}
对整数 \(m\ge 1\) 与路径长度 \(n\ge 0\)，记
\[
a_m(n):=\#\{\text{在 }\mathsf{Ind}(P_n)\text{ 上由 }c_n\text{ 生成的长度 }m\text{ 的循环}\}.
\]
对一般 \(x\)，同理记
\[
A_m(x):=\#\{\text{在纤维 }F_x\text{ 上由 }c_x\text{ 生成的长度 }m\text{ 的循环}\}.
\]
则对任意 \(m\ge 1\) 有精确公式
\[
A_m(x)=\sum_{\substack{m_1,\dots,m_r\ge 1\\ \mathrm{lcm}(m_1,\dots,m_r)=m}}
\Biggl(\ \prod_{j=1}^r a_{m_j}(\ell_j)\ \Biggr)\cdot \frac{\prod_{j=1}^r m_j}{m}.
\]
特别地，由定理 \ref{thm:pom-toggle-scan-cycle-type-closed} 的闭式 \(a_m(n)\) 可显式计算，从而 \(c_x\) 的全循环分解在分量数据 \((\ell_j)\) 下完全可计算。
\end{proposition}
\begin{proof}
固定一组分量循环 \(C_j\subset \mathsf{Ind}(P_{\ell_j})\)，其长度分别为 \(m_j\)。
则在 \(C_1\times\cdots\times C_r\) 上，\(c_x\) 同构于有限交换群
\(
\mathbb Z/m_1\mathbb Z\times\cdots\times \mathbb Z/m_r\mathbb Z
\)
上的平移
\[
(t_1,\dots,t_r)\longmapsto (t_1+1,\dots,t_r+1).
\]
该平移的步长元 \((1,\dots,1)\) 的阶为 \(\mathrm{lcm}(m_1,\dots,m_r)=m\)，故每条轨道长度为 \(m\)。
同时轨道条数等于
\[
\frac{|C_1\times\cdots\times C_r|}{m}=\frac{\prod_{j=1}^r m_j}{m}.
\]
对固定的 \((m_1,\dots,m_r)\)，从各分量中选择长度 \(m_j\) 的循环共有 \(\prod_{j=1}^r a_{m_j}(\ell_j)\) 种；
每一种选择贡献 \((\prod_j m_j)/m\) 条长度为 \(m\) 的循环。
对所有满足 \(\mathrm{lcm}(m_1,\dots,m_r)=m\) 的 \((m_1,\dots,m_r)\) 求和即得公式。证毕。
\end{proof}

\begin{theorem}[ghost 坐标的乘法化：\(\mathrm{lcm}\)-卷积的除子和线性化]\label{thm:pom-toggle-scan-cycle-ghost-multiplicativity}
沿用命题 \ref{prop:pom-toggle-scan-cycle-convolution} 的记号。
定义“循环质量”算术函数
\[
f_\ell(m):=m\,a_m(\ell),\qquad
f_x(m):=m\,A_m(x),
\]
并定义 \(\mathrm{lcm}\)-卷积
\[
(f\star_{\mathrm{lcm}}h)(n):=\sum_{\mathrm{lcm}(a,b)=n} f(a)\,h(b),
\]
以及 ghost 变换（除子和）
\[
g_\ell(n):=\sum_{d\mid n} f_\ell(d),\qquad
g_x(n):=\sum_{d\mid n} f_x(d).
\]
则：
\begin{enumerate}
  \item \(f_x\) 是 \((f_{\ell_1},\dots,f_{\ell_r})\) 的迭代 \(\mathrm{lcm}\)-卷积：
  \[
  f_x=f_{\ell_1}\star_{\mathrm{lcm}}\cdots\star_{\mathrm{lcm}} f_{\ell_r}.
  \]
  \item ghost 变换将 \(\mathrm{lcm}\)-卷积线性化为点乘：对任意 \(n\ge 1\)，
  \[
  \boxed{\;
  g_x(n)=\prod_{j=1}^r g_{\ell_j}(n)\; }.
  \]
\end{enumerate}
\end{theorem}
\begin{proof}
由命题 \ref{prop:pom-toggle-scan-cycle-convolution} 的闭式，两端同乘 \(m\) 得
\[
f_x(m)
=\sum_{\substack{m_1,\dots,m_r\ge 1\\ \mathrm{lcm}(m_1,\dots,m_r)=m}}
\prod_{j=1}^r f_{\ell_j}(m_j),
\]
这正是迭代 \(\mathrm{lcm}\)-卷积，得 (1)。
\par
对 (2)，对任意 \(n\ge 1\) 有
\[
g_x(n)
=\sum_{d\mid n} f_x(d)
=\sum_{\mathrm{lcm}(m_1,\dots,m_r)\mid n}\ \prod_{j=1}^r f_{\ell_j}(m_j).
\]
而 \(\mathrm{lcm}(m_1,\dots,m_r)\mid n\) 当且仅当对所有 \(j\) 都有 \(m_j\mid n\)。
因此
\[
g_x(n)
=\sum_{m_1\mid n}\cdots\sum_{m_r\mid n}\prod_{j=1}^r f_{\ell_j}(m_j)
=\prod_{j=1}^r\left(\sum_{m_j\mid n} f_{\ell_j}(m_j)\right)
=\prod_{j=1}^r g_{\ell_j}(n),
\]
证毕。
\end{proof}
\leanverified{paper\_pom\_toggle\_scan\_cycle\_ghost\_multiplicativity}
